\justifying
\NSFCBodyText

\subsubsubsection{研究内容}

\textbf{（1）三时钟观测机制下的运输事件风险辨识。} 面向振动、倾斜、电子封志、光感和包装图像，建立事件时间、设备采样时间与决策端到达时间的数据规范，显式建模状态相关采样和延迟交付。在可观测锚点和支持条件成立时估计运输关键状态、冲击/倾斜、非法开启、包装破损或变形的风险与时间区间；条件偏离时通过敏感性分析给出部分辨识区间或拒判，形成包含事件候选、证据完整度和不确定性的可靠证据状态。

\textbf{（2）截止期约束下的事件证据集合价值与安全资源决策。} 在 2G/3G 弱网、有限缓存、算力和能量条件下，研究候选证据相对于现有集合的截止期条件价值，刻画多模态证据的互补、冗余和迟到失效。以价值置信下界协调采样与数据压缩、离线缓存与状态保持、断点续传及边缘计算，并用端侧最低告警、关键原始证据保留、压缩保真和幂等续传约束稳定通信与可靠上传；以按期完整交付、通信字节、估算或实测能耗及续航代理量评价弱网和低功耗效果。

\textbf{（3）漂移环境下的仓储风险校准与条件性事件记忆。} 融合温湿度、安防视频、烟雾和门禁信息，分别刻画缓变暴露与突发风险，在仓库、季节和设备漂移下形成可校准、可拒判的实时监测状态，输出异常识别、分级预警和应急响应建议。仓储本地校准与拒判为必做主问题；货物级配对和统计功效达到进入条件后，再对冲击/倾斜—包装损伤、非法开启—货物安全等合理链条检验事件特异衰减与附加预测价值。

\subsubsubsection{研究目标}

总体目标是揭示弱网状态相关观测下可靠证据状态的形成、保全和后续风险使用规律，建立“可辨识—可保全—可校准”的跨境物流异常感知与边缘协同方法。具体目标为：形成能够区分未发生、未采到和未送达的三时钟风险辨识方法；形成以截止期集合条件价值为核心、受安全屏障约束的弱网资源策略；形成漂移条件下仓储多源风险校准、拒判和分级响应方法，并在配对门通过时检验条件性事件记忆。各目标分别以组平衡时间依赖 Brier 分数、截止期加权关键异常漏警风险、选择性风险—覆盖、证据链交付/资源指标及协议硬违例为验收口径，具体效应阈值由预实验和功效分析预先冻结。

\subsubsubsection{拟解决的关键科学问题}

\textbf{（1）状态相关缺失与三时钟错位下，已到达多模态证据在什么条件下仍能辨识运输事件风险？} 关键在于区分事件、采样和到达过程；条件不足时输出部分辨识区间或拒判，避免对未观测轨迹作过度自信的补全。

\textbf{（2）资源动作对事件证据集合的截止期条件价值如何度量？} 关键在于仅用当时可得信息刻画证据互补、冗余、交付概率和迟到失效，并使通信/计算/能量优化服从关键证据安全下界。

\textbf{（3）仓储多源风险在环境漂移下如何保持校准并形成可检验的分级响应？} 关键在于以仓储本地数据完成校准与拒判；只有实体配对条件具备时，才用摘要置换和时间外推检验运输事件记忆的附加预测价值，避免把普通相关性解释为因果作用。
